\section{MMDiT: que el texto y la imagen evolucionen juntos en espacios distintos}

En PixArt, los token de imagen se actualizan mediante cross-attention, y los token de texto suelen permanecer congelados. El MMDiT de Stable Diffusion 3 parte de que las dos modalidades están en espacios de representación distintos, por lo que requieren normalization y projections independientes, aunque pueden intercambiar información en la joint attention.

\subsection{Por qué no tratar las modalidades directamente como el mismo tipo de token}

El texto y la noisy image difieren en su estadística, en la longitud de su secuencia y en su semántica. Compartir por completo los parámetros podría obligar a una modalidad a heredar los sesgos de la otra; separarlos por completo impide la interacción. MMDiT opta por «proyecciones independientes, atención conjunta». Lo que esta sección busca distinguir es que compartir e interactuar no son lo mismo: las dos modalidades pueden conservar cada una su propia normalization/QKV y, aun así, intercambiar información en el mismo attention map; las cuatro diapositivas siguientes desarrollan paso a paso precisamente esa frontera.

\lecturefigure{slide-45.jpg}{Motivación de MMDiT: los espacios de embedding del texto y de la imagen son distintos, y la cross-attention unidireccional introduce un sesgo}{45}

La atención cruzada unidireccional solo actualiza el flujo de imagen; la representación del texto no cambia según el estado visual actual. MMDiT busca que ambos lados evolucionen juntos, conservando al mismo tiempo cada uno su propia parametrización. Esta motivación es razonable, pero no constituye una demostración teórica de que «separar siempre es mejor».

\lecturefigure{slide-46.jpg}{MMDiT usa AdaLN y QKV projections independientes para cada modalidad}{46}

El texto y la imagen se calculan por separado:
\[
(Q_x,K_x,V_x)=f_x(X,c),\qquad
(Q_t,K_t,V_t)=f_t(T,c),
\]
donde $f_x,f_t$ incluyen cada una su propia normalization y su propia proyección. Duplicar los parámetros aumenta la capacidad expresiva, pero también la memoria, los FLOPs y la complejidad de implementación.

\lecturefigure{slide-47.jpg}{Block completo de MMDiT: dos rutas de modulación, dos rutas de proyección y atención conjunta sobre los token}{47}

Tras concatenar $Q=[Q_t;Q_x]$, $K=[K_t;K_x]$ y $V=[V_t;V_x]$, se calcula la joint attention:
\[
O=\operatorname{softmax}\!\left(\frac{QK^\top}{\sqrt d}\right)V.
\]
Después, la salida se divide por modalidad. Así coexisten las interacciones text-text, image-image y la cross-modal interaction, en lugar de apilar por separado una capa de cross-attention.

\lecturefigure{slide-48.jpg}{Ventajas de MMDiT: cada modalidad conserva parámetros exclusivos y, a la vez, interactúa globalmente mediante la concatenación}{48}

Las separate paths permiten que cada modalidad aprenda su propio scale, shift y projections; la concatenation, en cambio, permite que la atención seleccione la información de manera unificada. El riesgo es el desequilibrio en la proporción de token: una gran cantidad de image tokens puede opacar un texto corto, lo cual exige un diseño de normalization, mask o loss.

\begin{knowledgebox}{Diferencia entre joint attention y cross-attention}
La cross-attention suele actualizar solo el query stream; la joint attention coloca los token de ambas modalidades en el mismo mapa de atención, de modo que ambos lados pueden actualizarse. La primera ahorra más cómputo y tiene una interfaz clara; la segunda ofrece una interacción más simétrica, pero con un costo mayor.
\end{knowledgebox}

\subsection{Resultados, escalamiento y componentes complementarios que suelen pasarse por alto}

Para saber si MMDiT vale la pena hay que fijarse en la ablation, no en el diagrama de la arquitectura. La clase presenta curvas de loss/quality y resultados de escalamiento, y advierte que SD3 también depende de mejores componentes de sampling/training; no se puede atribuir toda la mejora al multimodal block. Por eso, esta sección examina por separado tres tipos de evidencia: la ablación controlada responde por el mecanismo local, las curvas de scaling responden por la escalabilidad, y los ejemplos del sistema completo solo dan cuenta del desempeño final de la recipe combinada.

\lecturefigure{slide-49.jpg}{Ablation de MMDiT: efecto de la parametrización exclusiva por modalidad en las curvas de entrenamiento y de calidad}{49}

Las curvas «sugieren con fuerza», pero no demuestran, que MMDiT sea eficaz. Conviene verificar si el baseline está igualado en número de parámetros, FLOPs y duración del entrenamiento; si MMDiT es más grande, parte de la ganancia podría provenir de la capacity. Lo ideal es comparar dos grupos de experimentos: compute-matched y parameter-matched.

\lecturefigure{slide-50.jpg}{Scaling de MMDiT: los modelos grandes siguen beneficiándose, pero necesitan QK-Norm para estabilizar el entrenamiento}{50}

Al aumentar el hidden size y el número de capas, los attention logits pueden explotar. QK-Norm normaliza query/key y controla la escala del producto punto. Es un componente de estabilidad del entrenamiento: no aumenta directamente la capacidad de representación, pero puede determinar si un modelo grande converge o no.

\lecturefigure{slide-51.jpg}{El desempeño completo de SD3 también depende de los componentes complementarios de entrenamiento/muestreo que representa el «hippo»}{51}

Esta diapositiva humorística usa un hipopótamo para recordar que, además del artículo sobre la arquitectura, hay componentes complementarios como Rectified Flow, los datos, los caption, el VAE y la guidance. Al leer cualquier narrativa del tipo «tal block produjo SD3», hay que buscar activamente los omitted components y evitar atribuir el resultado a una sola causa.

\teachervoice{00:49:49--00:57:31 y 01:03:04--01:04:45: el docente describe la justificación de MMDiT como la diferencia entre los espacios de las modalidades y su evolución conjunta, pero en la sesión de preguntas reconoce que se trata de una hand-wavy intuition, no de una optimality proof.}

\begin{warningbox}{Explicar la arquitectura no es una descomposición causal}
Una curva de entrenamiento está influida a la vez por el número de parámetros, el optimizer, los datos, el loss weighting y la estabilidad numérica. Sin experimentos igualados, no puede afirmarse que «las separate QKV por sí solas causan toda la mejora».
\end{warningbox}

\subsection{Las diversas soluciones de compromiso de MMDiT}

El MMDiT completo mantiene dos juegos de projections en cada capa, lo cual es muy costoso. Los modelos reales pueden usar el multimodal block solo en algunas capas y el DiT ordinario en las demás; también pueden actualizar las modalidades de forma alternada, o bien unirlas hasta la etapa final. Esta sección pasa de la etiqueta binaria «usa o no MMDiT» al block schedule: hay que comparar en qué capas ocurre la unión, cuánto dura, cuánto se conserva de los parámetros de las dos rutas, y cómo estas decisiones modifican la Pareto frontier entre calidad y costo.

\lecturefigure{slide-52.jpg}{Variantes de MMDiT: combinación de multimodal blocks con DiT blocks ordinarios}{52}

Si solo hay $k$ capas MMDiT y las $L-k$ restantes son capas ordinarias, puede ajustarse el equilibrio entre la capacidad de interacción y los FLOPs. Modelos como Flux adoptan un diseño mixto. La pregunta central es si rinde más la fusión temprana, intermedia o tardía, y qué capas necesitan realmente parámetros de doble ruta.

\lecturefigure{slide-53.jpg}{Otra variante: primero se actualiza cada modalidad por separado y después las representaciones pasan a una etapa conjunta}{53}

En la figura, las modality A/B se procesan primero de forma independiente y luego se intercambian o se concatenan. Un diseño por etapas puede reducir la longitud en token de la joint attention, pero puede retrasar la alineación entre modalidades. Equivale a una elección continua entre early/mid/late fusion.

\lecturefigure{slide-54.jpg}{Modelos como Lumina-Image 2.0 muestran distintas combinaciones de multimodal blocks}{54}

La variedad de variantes muestra que MMDiT no es una receta fija, sino un espacio de diseño de «parámetros exclusivos por modalidad + cuándo unirlas». Al comparar artículos, conviene escribir el block schedule en una tabla, en lugar de solo marcar si se usa o no MMDiT.

\begin{lstlisting}[language=Python,caption={MMDiT block: QKV independientes, attention conjunta y división posterior por modalidad}]
def mmdit_block(text_tokens, image_tokens, condition):
    q_t, k_t, v_t = text_projection(text_tokens, condition)
    q_x, k_x, v_x = image_projection(image_tokens, condition)
    q = concat(q_t, q_x); k = concat(k_t, k_x); v = concat(v_t, v_x)
    output = attention(q, k, v)
    text_out, image_out = split_by_modality(output)
    return text_tokens + text_out, image_tokens + image_out
\end{lstlisting}

\subsection{Resumen de la sección}

MMDiT hace que el texto y la imagen usen AdaLN/QKV independientes y, al mismo tiempo, evolucionen juntos en la joint attention. Esto puede atenuar el conflicto entre espacios de representación, pero también aumenta notablemente el costo. Los sistemas reales recurren a blocks mixtos, a la fusión por etapas y a QK-Norm para equilibrar calidad, estabilidad y FLOPs.
